% =====================================================================
%  Section 5: Evaluation  --  paste into Overleaf
%  Requires in your preamble:
%     \usepackage{booktabs}
%     \usepackage{xspace}
%     \usepackage{amssymb}          % for \checkmark / \times if desired
%     \newcommand{\tool}{\textsc{SpecRL}\xspace}   % <-- rename to your system name
%     \newcommand{\TODO}[1]{\textbf{\color{red}[TODO: #1]}} % fill in real numbers
%  (If you use the todonotes package instead, drop the \TODO definition above.)
% =====================================================================

\section{Evaluation}
\label{sec:eval}

We evaluate \tool along two axes.
Section~\ref{sec:eval:config} describes the hardware and software platform.
Section~\ref{sec:eval:security} then presents the security evaluation:
we first show that \tool rediscovers known transient-execution attacks
on our target CPU (Section~\ref{sec:eval:repro}), validating that the
reinforcement-learning fuzzer drives the hardware into the intended
speculative regimes.

\subsection{Hardware and Software Configuration}
\label{sec:eval:config}

All experiments run on a desktop AMD~Zen~4 platform.
Table~\ref{tab:platform} summarises the hardware and software stack, and
Table~\ref{tab:mitigations} lists the CPU's speculative-execution
vulnerability status as reported by the Linux kernel
(\texttt{/sys/devices/system/cpu/vulnerabilities/}).
The kernel marks this part as \emph{Vulnerable} to Spectre-v1, Spectre-v2,
Speculative Store Bypass (v4), and SRSO, and \emph{Not affected} by
Meltdown and MDS-class leakage; this delimits the attack classes that a
correct fuzzer should be able to surface on this machine.

\begin{table}[t]
\centering
\caption{Evaluation platform.}
\label{tab:platform}
\small
\begin{tabular}{ll}
\toprule
\textbf{Component} & \textbf{Configuration} \\
\midrule
CPU            & AMD Ryzen 7 8700F (Zen~4, family \texttt{0x19}) \\
Cores/Threads  & 8 cores / 16 threads (SMT disabled) \\
Microcode      & \texttt{0xa705208} \\
L1d / L1i      & 32\,KiB / 32\,KiB per core \\
L2 / L3        & 1\,MiB per core / 16\,MiB shared \\
\midrule
OS             & Ubuntu 24.04.4 LTS \\
Kernel         & Linux 6.8.0-136-generic \\
Runtime        & Python 3.10, Unicorn 1.0.3 \\
RL stack       & Ray/RLlib 2.55 (PPO), PyTorch 2.11 \\
Fuzzer         & \tool, built on Revizor~\cite{revizor} \\
\bottomrule
\end{tabular}
\end{table}

\begin{table}[t]
\centering
\caption{Kernel-reported vulnerability status of the target CPU.}
\label{tab:mitigations}
\small
\begin{tabular}{ll}
\toprule
\textbf{Class} & \textbf{Status (kernel report)} \\
\midrule
Spectre-v1 (BCB)             & Vulnerable \\
Spectre-v2 (BTI)             & Vulnerable (IBPB / STIBP disabled) \\
Spec.\ Store Bypass (v4)     & Vulnerable \\
SRSO                         & Vulnerable \\
Meltdown                     & Not affected \\
MMIO Stale Data / MDS        & Not affected \\
\bottomrule
\end{tabular}
\end{table}

\paragraph{Measurement setup.}
To keep the hardware channel stable we disable simultaneous
multithreading (SMT), turn off the NMI watchdog, and pin the CPU
frequency governor to \texttt{performance}; the kernel-side executor
pins every measurement to a single physical core and serialises
repetitions. These steps suppress the cross-thread and
frequency-scaling noise that otherwise dominates the false-positive
rate. Hardware traces (\emph{htraces}) are collected with a
Prime+Probe monitor over the 64 sets of the L1 data cache, yielding a
per-input cache-set occupancy vector. Contract traces (\emph{ctraces})
are produced by a Unicorn-based emulator that implements the leakage
contract under test. For every candidate program each input is measured
with warm-up rounds followed by an adaptive sampling schedule that
escalates the sample size (up to $5000$) and re-primes when a signal is
borderline, so that a reported violation reflects a statistically stable
htrace difference rather than a transient measurement artifact.

\subsection{Security Evaluation}
\label{sec:eval:security}

\tool casts test-program generation as a reinforcement-learning
problem. A PPO agent emits an instruction sequence from a
Unicorn-modellable action space; for each generated program the
emulator computes the ctrace under a target leakage
contract while the executor measures the htrace on the real CPU.
Following the leakage-contract framework~\cite{contracts,revizor}, a
\emph{contract violation} is a pair of inputs that share the same
ctrace---i.e.\ are indistinguishable under the contract---yet yield
different htraces---i.e.\ the hardware \emph{does} distinguish
them---witnessing leakage beyond what the contract permits.
Same-ctrace input classes are built cheaply by taint-guided input
\emph{boosting}: the emulator records which input bytes flow into the
ctrace, and boosting holds those bytes fixed while randomising the
rest, so every boosted input in a class provably shares the ctrace.
The agent's reward is driven by htrace divergence within a ctrace
class, steering generation toward violating programs.

\subsubsection{Reproducing Existing Attacks}
\label{sec:eval:repro}

\paragraph{Methodology.}
To reproduce a known transient-execution attack with \tool we
(i)~enable the microarchitectural primitive it relies on in the program
generator, (ii)~fuzz against the \emph{weakest} contract---one that
forbids the corresponding speculation---so that any hardware
speculation manifests as a violation, and (iii)~\emph{attribute} the
violation to a specific variant by identifying the minimal execution
clause that, once added to the contract, makes the violation disappear.
Step~(iii) converts a raw leak signal into a mechanistic attribution and
rules out unrelated sources of speculation. Table~\ref{tab:repro}
summarises the two reproduced attacks.

\paragraph{Spectre-V1 (Bounds-Check Bypass).}
We configure the observation clause to \texttt{loads+stores+pc} and the
execution clause to \texttt{no\_speculation} (the sequential, non-speculative
contract), and enable conditional branches (\texttt{BASE-COND\_BR}) in the
instruction pool. \tool synthesises programs in which a conditional branch
guards a subsequent, address-dependent memory access; under branch
misprediction the guarded access executes transiently and imprints a
secret-dependent line into the cache, which the Prime+Probe monitor
observes. Because the sequential contract predicts no such access, the two
boosted inputs share a ctrace but differ in htrace, and \tool reports a
violation after \TODO{time / \#test cases to first V1 violation}.
Adding the \texttt{cond} clause---which models conditional-branch
misprediction---to the contract eliminates the violation, attributing the
leak to bounds-check-bypass speculation.

\paragraph{Spectre-V4 (Speculative Store Bypass).}
We set the observation clause to \texttt{ct} and the execution clause to
\texttt{seq}, and \emph{disable} the hardware store-bypass mitigation by
clearing the SSBD control (\texttt{x86\_executor\_enable\_ssbp\_patch =
false}). \tool synthesises store/load pairs that target aliasing sandbox
addresses; the younger load speculatively bypasses the older, unresolved
store and forwards stale memory, which then leaks through a dependent
access. The sequential contract forbids store-to-load bypass, so the
resulting htrace difference is flagged as a violation after
\TODO{time / \#test cases to first V4 violation}. We confirm the
attribution with two independent controls: (a)~adding the \texttt{bpas}
(speculative store-bypass) clause to the contract eliminates the
violation, and (b)~re-enabling the hardware SSBD mitigation
(setting the \texttt{SPEC\_CTRL.SSBD} bit) also eliminates it. Together
these establish that the leak is genuine speculative store bypass rather
than an emulator-modelling artifact.

\begin{table*}[t]
\centering
\caption{Reproduction of known transient-execution attacks with \tool
on AMD Zen~4. Each attack is fuzzed against the weakest (non-speculative)
contract; the ``Attribution'' column gives the minimal execution clause
that eliminates the violation.}
\label{tab:repro}
\small
\begin{tabular}{lllll l}
\toprule
\textbf{Variant} & \textbf{Primitive} & \textbf{Obs.\ clause} & \textbf{Exec.\ clause} & \textbf{Key config} & \textbf{Detected / Attribution} \\
\midrule
Spectre-V1 (BCB) & Conditional branch      & \texttt{loads+stores+pc} & \texttt{no\_speculation} & \texttt{COND\_BR} enabled & \checkmark\ (\TODO{}) / \texttt{cond} \\
Spectre-V4 (SSB) & Store-to-load bypass    & \texttt{ct}              & \texttt{seq}             & SSBD off                  & \checkmark\ (\TODO{}) / \texttt{bpas} \\
\bottomrule
\end{tabular}
\end{table*}

\paragraph{Summary.}
On a CPU the kernel reports as vulnerable to both variants, \tool
independently rediscovers Spectre-V1 and Spectre-V4 from a generic
action space, without any variant-specific templates, and attributes
each leak to its speculation primitive via the minimal explaining
contract. This validates that the RL-driven generator reaches the
speculative regimes of interest and that the contract-based oracle
correctly separates them.
